\documentclass[10pt,twocolumn,letterpaper]{article}

\usepackage[margin=0.75in]{geometry}
\usepackage{amsmath,amssymb,amsfonts,bm}
\usepackage{booktabs}
\usepackage{graphicx}
\usepackage{microtype}
\usepackage{hyperref}
\usepackage{xcolor}
\usepackage{cite}
\usepackage{abstract}

\hypersetup{
    colorlinks=true,
    linkcolor=blue,
    citecolor=blue,
    urlcolor=blue
}

\renewcommand{\abstracttextfont}{\small}

\title{\textbf{Causal Tracing and Circuit Discovery in Small Open-Weight Transformers: Isolating Factual Retrieval and Suppression Mechanisms}}

\author{
  \textbf{Vedansh Kumar} \\
  \textit{Independent Researcher} \\
  \texttt{vedanshk10@gmail.com}
}

\date{\today}

\begin{document}

\maketitle

\begin{abstract}
Understanding how autoregressive language models store, retrieve, and route factual knowledge remains a central challenge in mechanistic interpretability. In this work, we reverse-engineer the factual association circuit in \texttt{gpt2-small} ($124\text{M}$ parameters) across a structurally balanced, single-token counterfactual dataset. Using Logit Lens and Direct Logit Attribution (DLA), we uncover a pronounced late-layer accumulation dynamic where factual confidence peaks at Layer 10 ($\Delta \mathcal{L} = +23.8$) before being damped at the output stage. 

Through activation patching across $144$ attention heads, residual streams, and MLP blocks, we isolate the spatial and temporal mechanics of factual recall: the subject entity is maintained in the residual stream at the subject token position through Layer 7, undergoes causal transfer across Layers 8--9, and is projected into the vocabulary space at the final token position. 

We identify a sparse writing core consisting of just three attention heads (\texttt{L9H8}, \texttt{L10H0}, \texttt{L8H11}). Joint mean-ablation of this three-head subcircuit causes an $87.46\%$ collapse in logit difference ($\Delta \mathcal{L} = +10.46 \rightarrow +1.31$) and degrades task accuracy from $100.0\%$ to $28.0\%$, whereas non-circuit control heads exhibit zero degradation. Finally, we discover a dedicated negative damping circuit dominated by head \texttt{L10H7}, whose targeted ablation disinhibits model confidence by $+14.56\%$.
\end{abstract}

\section{Introduction}
Autoregressive Transformer language models exhibit high empirical competence on factual retrieval tasks, yet their internal representations remain opaque. Mechanistic interpretability seeks to demystify these architectures by decomposing complex neural computations into discrete, human-interpretable subgraphs termed \textit{circuits}~\cite{elhasy2021transformer, wang2022interpretability}.

While prior foundational studies have investigated entity association in large models~\cite{meng2022locating} or syntactic behaviors like Indirect Object Identification (IOI)~\cite{wang2022interpretability}, the fine-grained causal interactions between information movers, direct writing heads, and negative damping heads in constrained parameter regimes remain actively investigated.

In this paper, we present an end-to-end causal discovery study on \texttt{gpt2-small}. We formalize the residual stream as an additive linear communication channel, trace the factual representation trajectory via activation patching, and validate circuit necessity and sufficiency through knockout mean-ablation experiments.

\section{Theoretical Framework \& Methodology}

\subsection{The Residual Stream as a Linear Accumulator}
In a decoder-only transformer with $L$ layers, the residual stream hidden state $x_L[-1] \in \mathbb{R}^{d_{\text{model}}}$ at the final token index $-1$ is an additive sum over embeddings, attention outputs, and MLP projections:
\begin{equation}
x_L[-1] = x_0[-1] + \sum_{l=0}^{L-1} \left( \text{Attn}_l[-1] + \text{MLP}_l[-1] \right)
\end{equation}
Logits $\mathcal{L} \in \mathbb{R}^{|V|}$ are computed via unembedding matrix $\mathbf{W}_U$:
\begin{equation}
\mathcal{L} = \text{LayerNorm}(x_L[-1]) \, \mathbf{W}_U
\end{equation}

\subsection{The Evaluation Metric: Logit Difference}
To avoid non-linear softmax saturation artifacts, we measure performance using the \textit{Clean vs. Corrupted Logit Difference} ($\Delta \mathcal{L}$):
\begin{equation}
\Delta \mathcal{L}(x) = \mathcal{L}(x)_{y_{\text{clean}}} - \mathcal{L}(x)_{y_{\text{corr}}}
\end{equation}
where $y_{\text{clean}}$ is the true factual completion (e.g., \texttt{" Paris"}) and $y_{\text{corr}}$ is the counterfactual target (e.g., \texttt{" Berlin"}).

\subsection{Direct Logit Attribution (DLA)}
Given the linear nature of the unembedding projection, the directional unembedding difference vector is defined as $\mathbf{d}_U = \mathbf{W}_U[:, y_{\text{clean}}] - \mathbf{W}_U[:, y_{\text{corr}}]$. The direct contribution of any module output vector $\mathbf{v} \in \mathbb{R}^{d_{\text{model}}}$ is given by:
\begin{equation}
\text{DLA}(\mathbf{v}) = \mathbf{v} \cdot \left( \frac{\mathbf{d}_U}{\text{LN\_Scale}(x_L[-1])} \right)
\end{equation}

\subsection{Causal Interchange Activation Patching}
To measure the causal mediation of an internal component $C$ at token position $p$, we cache activations from clean prompts $\mathcal{D}_{\text{clean}}$ and inject them into forward passes of corrupted prompts $\mathcal{D}_{\text{corr}}$:
\begin{equation}
\text{Recovery}(C, p) = \frac{\Delta \mathcal{L}_{\text{patched}}(C, p) - \Delta \mathcal{L}_{\text{corrupted}}}{\Delta \mathcal{L}_{\text{clean}} - \Delta \mathcal{L}_{\text{corrupted}}}
\end{equation}
A recovery of $1.0$ indicates that component $C$ at position $p$ causally accounts for $100\%$ of the factual recall gap.

\subsection{Dataset Mean-Ablation}
To test circuit necessity without inducing out-of-distribution LayerNorm scaling artifacts, target head activations $\mathbf{a}_{l, h}[p]$ are replaced with their dataset mean:
\begin{equation}
\mathbf{a}_{l, h}^{(i)}[p] \leftarrow \bar{\mathbf{a}}_{l, h}[p] = \frac{1}{|\mathcal{D}|} \sum_{k=1}^{|\mathcal{D}|} \mathbf{a}_{l, h}^{(k)}[p]
\end{equation}

\section{Empirical Results}

\begin{figure*}[t]
\centering
\includegraphics[width=\textwidth]{figure1_dla_and_logit_lens.pdf}
\caption{\textbf{Direct Logit Attribution and Logit Lens Trajectory.} (A) Logit Lens across intermediate layers revealing a sharp transition at Layer 8 and a peak at Layer 10 ($\Delta \mathcal{L} = +23.8$). (B) Layer-by-layer decomposition into Attention and MLP components. (C) $12 \times 12$ head-level DLA heatmap highlighting candidate writing heads \texttt{L9H8}, \texttt{L10H0}, and \texttt{L8H11}.}
\label{fig:fig1}
\end{figure*}

\subsection{Dataset Construction \& Baseline Metrics}
We constructed an isomorphic benchmark of $N=50$ single-token factual counterfactual pairs using a calibrated two-shot schema. On this benchmark, unmodified \texttt{gpt2-small} achieves $100.0\%$ clean top-1 accuracy with a mean clean logit difference of $\Delta \mathcal{L} = +10.4602$.

\subsection{Residual Stream Dynamics \& Writer Sparsity}
As illustrated in Figure~\ref{fig:fig1}A, the Logit Lens indicates that factual signal remains minimal ($\Delta \mathcal{L} \approx +0.5$ to $+1.1$) through Layer 7. A steep upward inflection begins at Layer 8, climbing to a peak of $+23.8$ at Layer 10, before undergoing damping in Layers 11--12.

Direct Logit Attribution (Figure~\ref{fig:fig1}C) demonstrates that output writing is concentrated in a tiny fraction of attention heads:
\begin{itemize}
    \item \texttt{L9H8}: $\text{DLA} = +70.36$
    \item \texttt{L10H0}: $\text{DLA} = +53.88$
    \item \texttt{L8H11}: $\text{DLA} = +42.33$
\end{itemize}
Conversely, head \texttt{L10H7} acts as a strong negative suppressor ($\text{DLA} = -22.56$), directly opposing factual overconfidence.

\begin{figure*}[t]
\centering
\includegraphics[width=\textwidth]{figure2_activation_patching.pdf}
\caption{\textbf{Causal Tracing via Activation Patching.} (A) Residual stream patching across token positions and layers, showing the causal transfer from the subject position (Token 18) in early layers to the final position (Token 19) in late layers. (B) MLP patching heatmap. (C) Attention head output patching at the final prompt token position. (D) Head output patching at the subject token position.}
\label{fig:fig2}
\end{figure*}

\subsection{Causal Tracing of Information Flow}
The activation patching sweep (Figure~\ref{fig:fig2}A) isolates the causal path of factual retrieval:
\begin{enumerate}
    \item \textbf{Subject Representation (Layers 0--7):} Patching the residual stream at the subject token position alone achieves $\approx 1.00$ logit recovery.
    \item \textbf{Causal Transfer Locus (Layers 8--9):} The causal recovery transfers from the subject token index to the final prompt index.
    \item \textbf{Final Vocabulary Projection (Layers 10--11):} Causal weight is concentrated at the prompt-final position, where direct writing heads project the fact into logits.
\end{enumerate}

\begin{table}[h]
\centering
\small
\caption{\textbf{Empirical Knockout Ablation Results across 11 Experimental Conditions.} Logit diff drop (\%) measures necessity.}
\label{tab:ablation}
\begin{tabular}{@{}lcccc@{}}
\toprule
\textbf{Condition} & $\Delta \mathcal{L}$ & \textbf{Drop (\%)} & \textbf{Acc (\%)} & \textbf{Heads} \\ \midrule
Clean Baseline & 10.46 & 0.00 & 100.00 & 0 \\ \midrule
Ablate \texttt{L9H8} & 7.35 & 29.74 & 86.00 & 1 \\
Ablate \texttt{L10H0} & 7.64 & 26.96 & 100.00 & 1 \\
Ablate \texttt{L8H11} & 8.84 & 15.53 & 86.00 & 1 \\ \midrule
Ablate Top 2 Writers & 3.38 & 67.73 & 92.00 & 2 \\
\textbf{Ablate Core Circuit} & \textbf{1.31} & \textbf{87.46} & \textbf{28.00} & \textbf{3} \\ \midrule
Ablate \texttt{L10H7} & 11.45 & -9.44 & 100.00 & 1 \\
Ablate All Suppressors & 11.98 & -14.56 & 100.00 & 2 \\ \midrule
Control 1 (\texttt{L0H0}) & 10.46 & 0.04 & 100.00 & 1 \\
Control 2 (\texttt{L4H4}) & 10.46 & 0.03 & 100.00 & 1 \\
Control 3 (3 Random) & 10.46 & -0.04 & 100.00 & 3 \\ \bottomrule
\end{tabular}
\end{table}

\begin{figure}[t]
\centering
\includegraphics[width=\columnwidth]{figure3_ablation_experiments.pdf}
\caption{\textbf{Circuit Ablation Battery.} (A) Logit difference collapse under selective head mean-ablation. (B) Task top-1 accuracy degradation across conditions.}
\label{fig:fig3}
\end{figure}

\subsection{Circuit Isolation via Knockout Ablation}
To verify circuit necessity, we performed mean-ablation across 11 distinct conditions (Table~\ref{tab:ablation} and Figure~\ref{fig:fig3}).

Ablating individual writers produces moderate degradation ($15.53\%$ to $29.74\%$ logit drop). When all three core writing heads (\texttt{L9H8}, \texttt{L10H0}, \texttt{L8H11}) are ablated simultaneously, the circuit collapses: logit difference drops by **$87.46\%$** ($\Delta \mathcal{L} = 1.31$) and accuracy drops to **$28.0\%$**.

In contrast, mean-ablating random control heads (\texttt{L0H0}, \texttt{L4H4}, or three random heads simultaneously) induces $\le 0.04\%$ variance with zero accuracy loss, confirming the specificity of the isolated circuit.

Finally, ablating the negative suppressor heads (\texttt{L10H7} and \texttt{L11H10}) yields a **$+14.56\%$ increase in logit difference**, demonstrating the disinhibition of factual recall.

\section{Discussion \& Limitations}
Our results show that factual association in \texttt{gpt2-small} is not uniformly distributed, but rather governed by a sparse circuit:
\begin{itemize}
    \item \textbf{Upstream Routing:} Early-layer representations store the subject entity locally up to Layer 7.
    \item \textbf{Causal Transfer:} Mid-layer MLPs and attention blocks transfer entity attributes across token positions in Layers 8--9.
    \item \textbf{Output Writing \& Damping:} Three late-layer attention heads write the target factual token, while negative heads prevent over-saturation.
\end{itemize}

\paragraph{Limitations.} This study focused on single-token factual associations within \texttt{gpt2-small}. Future work will extend this causal tracing methodology to multi-token generation and larger architectures (e.g., Pythia-$2.8\text{B}$, Llama-3-$8\text{B}$).

\section{Conclusion}
We have localized and causally verified the factual retrieval subcircuit in \texttt{gpt2-small}. Through activation patching and mean-ablation, we demonstrated that factual recall is concentrated in an upstream subject-routing mechanism, a 3-head writing core whose knockout reduces logit difference by $87.46\%$, and a distinct negative suppressor system that calibrates final prediction confidence.

\bibliographystyle{IEEEtran}
\bibliography{references}

\end{document}